\documentclass[11pt]{article}
\usepackage[a4paper,margin=20mm]{geometry}
\usepackage{amsmath,amssymb,xcolor,graphicx}
\usepackage[hypcap=false]{caption}
\usepackage[hidelinks,bookmarksnumbered,bookmarksopen]{hyperref}
\hypersetup{pdftitle={Surface analysis — X-ray photoelectron spectroscopy (XPS)},pdfauthor={Lecturer: Dr L. Moreau},
  pdfcreator={KherveNote}}
\renewcommand\labelitemii{\textbullet}
\renewcommand\labelitemiii{\textbullet}
\renewcommand\labelitemiv{\textbullet}
\renewcommand\labelenumii{\arabic{enumi}.\arabic{enumii}.}
\renewcommand\labelenumiii{\arabic{enumi}.\arabic{enumii}.\arabic{enumiii}.}
\renewcommand\labelenumiv{\arabic{enumi}.\arabic{enumii}.\arabic{enumiii}.\arabic{enumiv}.}
\setlength{\parindent}{0pt}
\setlength{\parskip}{0.6em}
\definecolor{knoteaccent}{HTML}{1A6DD8}
\definecolor{knotemuted}{HTML}{6B6B6B}
\definecolor{knotekey}{HTML}{D96B00}
\newcommand\knotetime[1]{\leavevmode\llap{\color{knotemuted}\scriptsize #1\hspace{1em}}}
\newcommand\knotekey[2][]{\par\noindent#1\colorbox{knotekey!12}{\parbox{\dimexpr\linewidth-2\fboxsep}{%
  \textbf{\color{knotekey}$\star$ Key point.} #2}}\par}
\newcommand\knotequestion[2][]{\par\noindent#1{\color{knoteaccent}\textbf{?}}~\emph{#2}\par}
\newenvironment{knotetranscript}{\par\color{knotemuted}\small}{\par}
\newcommand\knotesummary[1]{\par\noindent\fcolorbox{knoteaccent}{knoteaccent!6}{%
  \parbox{\dimexpr\linewidth-2\fboxsep-2\fboxrule}{\textbf{Summary}\par #1}}\par}
\newenvironment{knotechunk}{}{}
\pagestyle{plain}
\begin{document}
\begin{knotechunk}
{\color{knoteaccent}\rule{\linewidth}{1.5pt}}\par
{\LARGE\bfseries Surface analysis — X-ray photoelectron spectroscopy (XPS)}\par
{\large Lecturer: Dr L. Moreau}\hfill{\color{knotemuted}09 October 2026 \textperiodcentered{} Materials building, LG11}\par
{\color{knoteaccent}\rule{\linewidth}{0.6pt}}\par
\knotesummary{XPS measures the binding energies of core electrons ejected by X-rays, giving the elements present in the top few nanometres and their chemical states. Peak positions shift with oxidation state; spin–orbit doublets have fixed area ratios; quantification needs a background and sensitivity factors.}

\section{Principle}

\knotetime{01:55}A photon of energy $h\nu$ ejects a core electron; we measure its kinetic energy: \[E_B = h\nu - E_K - \phi_{\mathrm{sp}}\] with $\phi_{\mathrm{sp}}$ the spectrometer work function.

\begin{itemize}\setlength{\itemsep}{0pt}\setlength{\parskip}{0pt}
  \item Al K\ensuremath{\alpha} (monochromated) 1486.6 eV; Mg K\ensuremath{\alpha} 1253.6 eV
  \item every element except H and He can be detected
\end{itemize}

\knotekey[\knotetime{03:10}]{Surface sensitive: photoelectrons travel only \textasciitilde{}1–3 nm (the inelastic mean free path \ensuremath{\lambda}) without losing energy, so \textasciitilde{}95 \% of the signal comes from within 3\ensuremath{\lambda}, i.e. the top \ensuremath{\approx} 5–10 nm.}
\end{knotechunk}

\begin{knotechunk}
\section{Chemical shifts}

\knotetime{10:25}The binding energy rises with the oxidation state (less electron density, less screening).

\begin{itemize}\setlength{\itemsep}{0pt}\setlength{\parskip}{0pt}
  \item Si 2p: Si\textsuperscript{0} \ensuremath{\approx} 99.3 eV, SiO\textsubscript{2} \ensuremath{\approx} 103.3 eV — about +4 eV
  \item C 1s: C–C \ensuremath{\approx} 284.8 eV, C–O \ensuremath{\approx} 286.5 eV, C=O \ensuremath{\approx} 288 eV, O–C=O \ensuremath{\approx} 289 eV
\end{itemize}

\subsection{Spin–orbit splitting}

\knotetime{12:05}Levels with $l > 0$ split into doublets with fixed area ratios:

\begin{enumerate}\setlength{\itemsep}{0pt}\setlength{\parskip}{0pt}
  \item p: $p_{1/2} : p_{3/2} = 1 : 2$
  \item d: $d_{3/2} : d_{5/2} = 2 : 3$
  \item f: $f_{5/2} : f_{7/2} = 3 : 4$
\end{enumerate}
\begin{itemize}\setlength{\itemsep}{0pt}\setlength{\parskip}{0pt}
  \item Au 4f\textsubscript{7}/\textsubscript{2} = 84.0 eV, splitting 3.67 eV — a common calibration
\end{itemize}
\end{knotechunk}

\begin{knotechunk}
\section{Peak fitting and quantification}

\begin{itemize}\setlength{\itemsep}{0pt}\setlength{\parskip}{0pt}
  \item subtract a background first: Shirley (iterative, step-like), Tougaard (physical loss function) or linear
  \item fit components with Voigt / pseudo-Voigt shapes; constrain doublet ratios and splittings
  \item atomic fraction $x_i = \dfrac{I_i/S_i}{\sum_j I_j/S_j}$ with relative sensitivity factors $S$
\end{itemize}

\knotekey[\knotetime{23:40}]{Accuracy is typically 10–20 \% relative — report it with the fitting choices.}
\end{knotechunk}

\begin{knotechunk}
\section{Practical issues}

\begin{itemize}\setlength{\itemsep}{0pt}\setlength{\parskip}{0pt}
  \item insulators charge positively: use a flood gun and reference the scale
  \item adventitious carbon C 1s at 284.8 eV is the usual reference, but it is debated
  \item sputtering for depth profiles can reduce oxides — beware of beam damage
\end{itemize}

\knotequestion[\knotetime{34:40}]{Is the adventitious carbon reference reliable on samples that were heated or sputtered?}
\end{knotechunk}

\begin{knotechunk}
\section{What was said}
\begin{knotetranscript}
\knotetime{09:30:50}XPS is our workhorse for surfaces, so it's worth understanding where every number comes from.

\knotetime{09:35:30}We use monochromated aluminium K alpha at fourteen eighty-six point six electronvolts.

\knotetime{09:38:10}The electrons we count have only come from the top few nanometres, deeper ones lose energy on the way out.

\knotetime{09:42:40}Oxidise silicon and the two p peak moves up by about four electronvolts, that's the chemical shift.

\knotetime{09:47:20}For p, d and f levels you always get a doublet, and the area ratio is fixed by the degeneracy.

\knotetime{09:50:30}Before fitting anything you need a background. Shirley is the common choice, Tougaard is more physical.

\knotetime{09:57:00}Quantification divides each area by its sensitivity factor; don't expect better than ten or twenty percent.

\knotetime{10:01:40}Insulating samples charge up, so we use the flood gun and reference to adventitious carbon.

\knotetime{10:05:10}There's an ongoing debate on whether two eight four point eight is a good reference at all.
\end{knotetranscript}
\end{knotechunk}
\end{document}
